\documentclass{article}
\usepackage{graphicx} % Required for inserting images
\usepackage{amsmath}
\usepackage{amssymb}
\usepackage{amsthm}
\usepackage{hyperref}

\title{Hausdorff Spaces}
\author{Aaron Honculada}
\date{February 2025}

\begin{document}

\maketitle

A \textbf{Hausdorff space} (also called a \textbf{T2 space}) is a topological space in which, for any two distinct points,
there exists disjoint open neighborhoods around each point. This property ensures that points in the space are ``separated"
in a certain sense, which is an important condition in many areas of topology and analysis.

\section{Formal Definition}
\hypertarget{hausdorff-definition}{}
A topological space $X$ is \textbf{Hausdorff} if, for any two distinct points $x, y \in X$, there exist open
sets $U$ and $V$ such that:
\begin{enumerate}
    \item $x \in U$
    \item $y \in V$
    \item $U \cap V = \emptyset$ (i.e. $U$ and $V$ are disjoint).
\end{enumerate}
In other words, you can find disjoint open neighborhoods for any two distinct points in the space, meaning that
the space is ``separated" enough to allow for a clear distinction between points.

\section{Why is this important?}
\begin{itemize}
    \item The Hausdorff condition is crucial for ensuring the uniqueness of limits. In Hausdorff spaces, sequences
    (or nets) of points that converge to a limit do so in a well-defined way, and the limit of a converging
    sequence is unique.
    \item Many important spaces in mathematics, such as metric spaces, Euclidean spaces, and manifolds, are Hausdorff
    because this property is desireable for the well-behavedness of convergence and continuity.
\end{itemize}


\section{Example of Hausdorff Space}
\begin{itemize}
    \item Euclidean space $\mathbb{R}^n$ (with the standard topology) is Hausdorff because
    for any two distinct points we can always find disjoint open neighborhoods around them. For example,
    in $\mathbb{R}^2$ you can choose two small balls centered at the points that do not overlap.
\end{itemize}

\section{Example of Non-Hausdorff Space}
\begin{itemize}
    \item The \textbf{co-finite topology} on an infinite set $X$, where the open sets are either $X$ itself or
    subsets of $X$ with finite complements is \textbf{not Hausdorff}. In this topology, for any two distinct points
    $x$ and $y$, any open sets containing them both must be large (since their complements are finite),
    so it is impossible to separate them with disjoint open sets.
\end{itemize}




\end{document}